% !TeX program = xelatex
\documentclass[a4paper,11pt]{article}

% ---------- Font & ngôn ngữ (Tiếng Việt qua XeLaTeX) ----------
\usepackage{fontspec}
\usepackage{polyglossia}
\setmainlanguage{vietnamese}
\setmainfont{Times New Roman}

% ---------- Trang & bố cục ----------
\usepackage[a4paper,margin=2.3cm]{geometry}
\usepackage{setspace}
\onehalfspacing
\usepackage{parskip}
\usepackage{titlesec}
\titlespacing*{\section}{0pt}{14pt}{6pt}
\titlespacing*{\subsection}{0pt}{10pt}{4pt}

% ---------- Bảng, màu, code ----------
\usepackage{xcolor}
\usepackage{array}
\usepackage{enumitem}
\usepackage{hyperref}
\usepackage{fancyhdr}
\usepackage{listings}

\definecolor{codebg}{RGB}{245,245,245}
\definecolor{codeframe}{RGB}{200,200,200}
\definecolor{milestone}{RGB}{31,73,125}
\definecolor{warnred}{RGB}{176,42,42}

\lstset{
  backgroundcolor=\color{codebg},
  basicstyle=\ttfamily\small,
  breaklines=true,
  frame=single,
  rulecolor=\color{codeframe},
  columns=fullflexible,
  showstringspaces=false,
  xleftmargin=8pt,
  xrightmargin=8pt,
  framextopmargin=4pt,
  framexbottommargin=4pt,
}

\hypersetup{
  colorlinks=true,
  linkcolor=milestone,
  urlcolor=milestone,
  citecolor=milestone,
  pdftitle={Loi giai tham khao - Do an DevOps VIMARU News Portal},
}

\pagestyle{fancy}
\setlength{\headheight}{14.5pt}
\fancyhf{}
\rhead{Đồ án DevOps — VIMARU News Portal}
\lhead{Lời giải tham khảo (nội bộ)}
\cfoot{\thepage}
\renewcommand{\headrulewidth}{0.4pt}

\title{\textbf{Lời giải tham khảo — Mốc 2}\\[4pt]
\large Dockerfile \& Docker Compose\\[2pt]
\normalsize Đồ án DevOps: VIMARU News Portal}
\author{}
\date{}

\begin{document}
\maketitle
\thispagestyle{fancy}

\begin{center}
\fbox{\parbox{0.85\linewidth}{\centering
\textcolor{warnred}{\textbf{Tài liệu nội bộ — không phát cho sinh viên.}}\\
Đây là bản lời giải đầy đủ, đã build và chạy thử thành công, dùng để giảng viên đối chiếu khi chấm Mốc 2. Các file \texttt{Dockerfile}, \texttt{.dockerignore}, \texttt{nginx.conf} và \texttt{docker-compose.yml} tương ứng \textbf{đã bị loại khỏi Git} (xem \texttt{.gitignore}) trong mã nguồn phát cho sinh viên — sinh viên phải tự viết lại từ đầu.
}}
\end{center}

\section*{1. \texttt{backend/Dockerfile}}

\begin{lstlisting}[language=bash]
FROM node:20-alpine
WORKDIR /app

COPY package*.json ./
RUN npm install --production

COPY . .

EXPOSE 4000
CMD ["node", "src/app.js"]
\end{lstlisting}

\section*{2. \texttt{auth-service/Dockerfile}}

Giống hệt mẫu backend, chỉ khác cổng:

\begin{lstlisting}[language=bash]
FROM node:20-alpine
WORKDIR /app

COPY package*.json ./
RUN npm install --production

COPY . .

EXPOSE 5000
CMD ["node", "src/app.js"]
\end{lstlisting}

\section*{3. \texttt{frontend/Dockerfile} (multi-stage: build rồi serve bằng Nginx)}

\begin{lstlisting}[language=bash]
# Stage 1: build ung dung React (Vite)
FROM node:20-alpine AS build
WORKDIR /app

COPY package*.json ./
RUN npm install

COPY . .
ARG VITE_API_BASE_URL=/api
ARG VITE_AUTH_BASE_URL=/auth
ENV VITE_API_BASE_URL=$VITE_API_BASE_URL
ENV VITE_AUTH_BASE_URL=$VITE_AUTH_BASE_URL
RUN npm run build

# Stage 2: serve file tinh bang Nginx
FROM nginx:alpine
COPY --from=build /app/dist /usr/share/nginx/html
COPY nginx.conf /etc/nginx/conf.d/default.conf

EXPOSE 80
\end{lstlisting}

Kèm theo \texttt{frontend/nginx.conf}:

\begin{lstlisting}[language=bash]
server {
    listen 80;
    server_name _;
    root /usr/share/nginx/html;
    index index.html;

    location / {
        try_files $uri $uri/ /index.html;
    }
}
\end{lstlisting}

\section*{4. \texttt{docker-compose.yml} (gốc repo, chạy đồng thời cả 3 service)}

\begin{lstlisting}[language=bash]
services:
  auth:
    build: ./auth-service
    container_name: vimaru-auth
    ports:
      - "5000:5000"
    environment:
      - PORT=5000
      - ADMIN_PASSWORD=admin123
      - JWT_SECRET=${JWT_SECRET:-vimaru-jwt-secret-change-me}
      - JWT_EXPIRES_IN=2h
    restart: unless-stopped

  backend:
    build: ./backend
    container_name: vimaru-backend
    ports:
      - "4000:4000"
    environment:
      - PORT=4000
      - JWT_SECRET=${JWT_SECRET:-vimaru-jwt-secret-change-me}
    volumes:
      - ./backend/data:/app/data
    depends_on:
      - auth
    restart: unless-stopped

  frontend:
    build:
      context: ./frontend
      args:
        VITE_API_BASE_URL: ${VITE_API_BASE_URL:-http://localhost:4000/api}
        VITE_AUTH_BASE_URL: ${VITE_AUTH_BASE_URL:-http://localhost:5000/api}
    container_name: vimaru-frontend
    ports:
      - "${FRONTEND_PORT:-3000}:80"
    depends_on:
      - backend
      - auth
    restart: unless-stopped
\end{lstlisting}

\section*{5. Điểm mấu chốt khi chấm}

\begin{itemize}[itemsep=3pt]
  \item \texttt{JWT\_SECRET} của \texttt{backend} và \texttt{auth} \textbf{phải giống nhau} — đây là lỗi phổ biến nhất khiến đăng nhập quản trị "chạy được nhưng thao tác thêm/sửa/xóa tin tức luôn báo 401".
  \item \texttt{frontend} cần build với \textbf{cả hai} biến \texttt{VITE\_API\_BASE\_URL} và \texttt{VITE\_AUTH\_BASE\_URL} — thiếu biến thứ hai thì trang đăng nhập quản trị sẽ gọi nhầm về địa chỉ mặc định.
  \item \texttt{docker-compose.yml} của sinh viên không bắt buộc phải giống hệt bản trên (tên service, thứ tự khai báo... có thể khác) — chỉ cần build được cả 3 image, `depends\_on` hợp lý, và 3 service giao tiếp đúng như mô tả trong tài liệu \texttt{nhiem-vu-sinh-vien.tex}.
\end{itemize}

\end{document}
